\chapter{Fecal Chymotrypsin}

\section*{What Is This Test?}
A fecal chymotrypsin test measures chymotrypsin activity in stool. Chymotrypsin is a pancreatic digestive enzyme secreted into the duodenum as chymotrypsinogen, then activated by trypsin and intestinal enteropeptidase \cite{layer1999pancreatic}. Unlike serum tests, fecal chymotrypsin directly reflects active enzyme entering the intestinal lumen and is therefore a marker of exocrine pancreatic function. The test is used to detect pancreatic exocrine insufficiency (PEI) in chronic pancreatitis, cystic fibrosis, and pancreatic cancer, and may be used to monitor pancreatic enzyme replacement therapy (PERT) adherence \cite{bowles2003role}.

\section*{Alternative Names}
Fecal Chymotrypsin; Stool Chymotrypsin; Stool Chymotrypsin Activity

\section*{Principle}
Stool is diluted and mixed with a synthetic substrate (casein or N-benzoyl-L-tyrosine ethyl ester). Color development is proportional to chymotrypsin activity. Photometric or turbidimetric assay at 405 nm. Unlike fecal elastase, chymotrypsin is degraded by intestinal bacteria, so values may be lower with prolonged transit time.

\section*{History}
Fecal chymotrypsin was introduced in the 1960s as a seminal pancreatic function test. By the 1990s, fecal elastase-1 largely replaced it because elastase is more stable in stool, less affected by intestinal degradation, and unaffected by PERT \cite{sellers2021fecal}. Fecal chymotrypsin remains useful, where elastase is not available, and specifically for monitoring PERT (chymotrypsin of porcine origin cross-reacts with assay).

\section*{Why This Test Is Ordered}
\begin{itemize}
  \item Diagnosis of pancreatic exocrine insufficiency in chronic pancreatitis \cite{dominguezmunoz2013diagnosis}
  \item Cystic fibrosis pancreatic sufficiency status
  \item Post-pancreatectomy or post-pancreatic surgery insufficiency assessment
  \item Pancreatic cancer with mass effect obstructing pancreatic duct
  \item Monitoring compliance with pancreatic enzyme replacement therapy (PERT)
\end{itemize}

\section*{Primary vs Secondary Test}
Secondary test; fecal elastase-1 is preferred where available because of better stability and specificity.

\section*{How This Test Is Performed}
A small stool sample (5--10 g) is collected in a special container without preservative, kept refrigerated, and assayed within 24--72 hours. Samples affected by water or urine are unacceptable. When monitoring PERT, samples taken while the patient is on therapy are interpreted bare in mind: chymotrypsin in stool will increase due to exogenous enzyme.

\section*{What Preparation Is Needed}
Stop pancreatic enzyme replacement therapy (PERT) 3--5 days before testing if measuring endogenous function. If PERT adherence is being monitored, continue it. Avoid barium or iodinated contrast studies within 72 hours.

\section*{Contraindications and Precautions}
False-low results: prolonged intestinal transit, antibiotic therapy suppressing proteolytic bacteria that activate the proenzyme, acid-suppressing medications, ileus. False-high: porcine PERT enzymes cross-react. Aspirin and barium interfere with the colorimetric assay.

\section*{Risks}
None; stool collection only.

\section*{What Happens After the Test}
Results available within 1--3 days. Low levels prompt evaluation of pancreatic structure (imaging) and consideration of PERT.

\section*{Understanding Results}

\subsection*{Below Normal}
\begin{itemize}
  \item Chronic pancreatitis with exocrine insufficiency
  \item Cystic fibrosis (late-stage pancreatic-insufficient CF)
  \item Pancreatic or ampullary cancer
  \item Pancreatic resection
  \item Severe celiac disease or Crohn disease with secondary pancreatic dysfunction
  \item Severe malnutrition
\end{itemize}

\subsection*{Above Normal}
Usually clinically unremarkable. High values if continuing PERT; not meaningful beyond compliance confirmation. Effective pancreatic function with no other clinical significance.

\section*{Normal Results}
\begin{itemize}
  \item \textbf{Fecal chymotrypsin:} $>$ 3 U/g stool (normal pancreatic exocrine function)
  \item \textbf{Moderate insufficiency:} 1.5--3 U/g stool
  \item \textbf{Severe insufficiency:} $<$ 1.5 U/g stool
  \item Interference: discontinue PERT 3--5 days before testing for endogenous function; chymotrypsin level reflects pancreatic enzyme secretion if exogenous replaced
\end{itemize}

\section*{Follow-up Tests}
\begin{itemize}
  \item \textbf{Fecal elastase-1:} More specific and stable marker of pancreatic exocrine insufficiency; preferred alternative
  \item \textbf{Secretin or cholecystokinin stimulation test:} Direct measurement of pancreatic secretion in duodenum (gold standard)
  \item \textbf{Abdominal CT or MRI:} For chronic pancreatitis structural changes (calcifications, duct dilatation)
  \item \textbf{MRCP or endoscopic ultrasound:} Detailed pancreatic duct imaging
  \item \textbf{Genetic testing (CFTR):} If cystic fibrosis suspected and not yet diagnosed
  \item \textbf{Oral pancreatic enzyme replacement therapy (PERT):} Confirm diagnosis before initiating
\end{itemize}

\section*{References}
\bibliographystyle{unsrtnat}
\bibliography{references}

\clearpage
\section*{Regulatory Status and Authorization Date}
Regulatory status applies to the specific assay, instrument, manufacturer, and intended use rather than to the test name alone. No single approval date can be assigned to this test category. For a commercial assay, verify the current product in the relevant regulator's database: the U.S. Food and Drug Administration (FDA) clearance, approval, or authorization; India's Central Drugs Standard Control Organisation (CDSCO) licence or permission; the European Union In Vitro Diagnostic Regulation (EU IVDR) CE certificate issued through the applicable conformity-assessment route; or the United Kingdom Medicines and Healthcare products Regulatory Agency (MHRA) registration and UKCA/CE status. Laboratory-developed tests may not have an individual FDA, CDSCO, EU IVDR, or MHRA approval date.
\clearpage
